\section*{NeurIPS Paper Checklist}

\begin{enumerate}

\item {\bf Claims}
    \item[] Question: Do the main claims made in the abstract and introduction accurately reflect the paper's contributions and scope?
    \item[] Answer: \answerYes{}
    \item[] Justification: The abstract and Section~\ref{sec:intro} state three contributions (SF-GSI, the GRAMLANG benchmark, and the billboard finding) that are each substantiated in Sections~\ref{sec:method}--\ref{sec:results}; quantitative claims (e.g., $89.7\%$ billboard, $0.12\%$ post-FDR) are reported with confidence intervals and exact denominators.

\item {\bf Limitations}
    \item[] Question: Does the paper discuss the limitations of the work performed by the authors?
    \item[] Answer: \answerYes{}
    \item[] Justification: A dedicated Limitations paragraph in Section~6 enumerates probe-class limits (linear/MLP heads on frozen embeddings), motif-call dependence on FIMO/JASPAR, and the dataset coverage gap (no developmental MPRA).

\item {\bf Theory assumptions and proofs}
    \item[] Question: For each theoretical result, does the paper provide the full set of assumptions and a complete (and correct) proof?
    \item[] Answer: \answerNA{}
    \item[] Justification: The paper is empirical; it introduces an estimator (SF-GSI, Eq.~\ref{eq:sfgsi}) but makes no theoretical claims about it.

\item {\bf Experimental result reproducibility}
    \item[] Question: Does the paper fully disclose all the information needed to reproduce the main experimental results?
    \item[] Answer: \answerYes{}
    \item[] Justification: Section~\ref{sec:setup} specifies models, datasets, motif caller and parameters, probe architecture/optimizer, and per-enhancer permutation count; appendices~\ref{app:probes}--\ref{app:sensitivity} document probe quality, null construction, and sensitivity analyses. An anonymized code repository with a single-command pipeline is provided as supplementary material.

\item {\bf Open access to data and code}
    \item[] Question: Does the paper provide open access to the data and code, with sufficient instructions to faithfully reproduce the main experimental results?
    \item[] Answer: \answerYes{}
    \item[] Justification: All code (probes, perturbation operators, SF-GSI, figure generation) is released under MIT at an anonymized URL provided in the supplementary material; the five MPRA datasets are public and cited; \texttt{environment.yml} pins all dependencies and \texttt{scripts/run\_full\_pipeline.py} reproduces every figure.

\item {\bf Experimental setting/details}
    \item[] Question: Does the paper specify all the training and test details necessary to understand the results?
    \item[] Answer: \answerYes{}
    \item[] Justification: Probe optimizer (Adam), weight decay ($10^{-3}$), training-set size ($5{,}000$ per dataset), motif-call $p$-value cutoff ($10^{-4}$), and permutation count ($N=30$) are stated in Sections~\ref{sec:method}--\ref{sec:setup}; full details and held-out splits are in Appendix~\ref{app:probes}.

\item {\bf Experiment statistical significance}
    \item[] Question: Does the paper report error bars suitably and correctly defined or other appropriate information about the statistical significance of the experiments?
    \item[] Answer: \answerYes{}
    \item[] Justification: Headline census reports a $95\%$ bootstrap confidence interval ($B=10{,}000$ enhancer-level resamples; Section~\ref{sec:census}); positive controls report $t$/$F$-test $p$-values and Cohen's $d$; per-enhancer SF-GSI uses Benjamini--Hochberg FDR.

\item {\bf Experiments compute resources}
    \item[] Question: Does the paper provide sufficient information on the computer resources needed to reproduce the experiments?
    \item[] Answer: \answerYes{}
    \item[] Justification: Section~\ref{sec:setup} and Appendix~D state a single A100 ($40$\,GB) GPU, total compute $<200$ GPU-hours, full-pipeline wall-clock $\sim 2$ days.

\item {\bf Code of ethics}
    \item[] Question: Does the research conducted in the paper conform, in every respect, with the NeurIPS Code of Ethics?
    \item[] Answer: \answerYes{}
    \item[] Justification: All experiments are in-silico on public MPRA data with no human subjects; the paper draws conservative conclusions (errs toward calling models grammar-insensitive) consistent with NeurIPS guidance on honest characterization of model capabilities.

\item {\bf Broader impacts}
    \item[] Question: Does the paper discuss both potential positive societal impacts and negative societal impacts of the work performed?
    \item[] Answer: \answerYes{}
    \item[] Justification: The Broader Impact paragraph in Section~6 notes that overestimating regulatory-grammar competence risks misleading downstream uses (synthetic enhancer design, variant interpretation, therapeutic target prioritization) and frames the paper's diagnostic as a mitigation.

\item {\bf Safeguards}
    \item[] Question: Does the paper describe safeguards that have been put in place for responsible release of data or models that have a high risk for misuse?
    \item[] Answer: \answerNA{}
    \item[] Justification: The paper releases analysis code and evaluation tooling only; no new generative model, no scraped dataset, and no high-misuse-risk artifact is released.

\item {\bf Licenses for existing assets}
    \item[] Question: Are the creators or original owners of assets used in the paper properly credited and are the license and terms of use explicitly mentioned and properly respected?
    \item[] Answer: \answerYes{}
    \item[] Justification: All five foundation models (DNABERT-2, NT, HyenaDNA, Enformer, PARM) and all five MPRA datasets are cited in Section~\ref{sec:setup} with their original publications; JASPAR2024 and FIMO are cited; the released analysis code uses MIT license, compatible with the licenses of the upstream tools.

\item {\bf New assets}
    \item[] Question: Are new assets introduced in the paper well documented and is the documentation provided alongside the assets?
    \item[] Answer: \answerYes{}
    \item[] Justification: The GRAMLANG benchmark and SF-GSI implementation are released with a \texttt{README.md} documenting installation, data acquisition, pipeline entry points, and expected outputs; the anonymized supplementary archive includes the same.

\item {\bf Crowdsourcing and research with human subjects}
    \item[] Question: For crowdsourcing experiments and research with human subjects, does the paper include the full text of instructions, screenshots, and compensation details?
    \item[] Answer: \answerNA{}
    \item[] Justification: The paper involves no crowdsourcing and no human subjects.

\item {\bf Institutional review board (IRB) approvals or equivalent for research with human subjects}
    \item[] Question: Does the paper describe potential risks to participants, whether risks were disclosed, and whether IRB approval was obtained?
    \item[] Answer: \answerNA{}
    \item[] Justification: The paper involves no human subjects.

\item {\bf Declaration of LLM usage}
    \item[] Question: Does the paper describe the usage of LLMs if it is an important, original, or non-standard component of the core methods?
    \item[] Answer: \answerNA{}
    \item[] Justification: LLMs were not used as part of the core methodology. The core method (SF-GSI) is a spacer-factored permutation test; the foundation models being probed are DNA models, not LLMs in the natural-language sense.

\end{enumerate}
